\section*{Overview}

Polynomial interpolation appears in contest problems whenever a sequence is secretly generated by a low-degree
polynomial. If you know enough sample values, you can evaluate the polynomial at new points without reconstructing every
coefficient explicitly.

\section*{When to Use It}

Use it when:

\begin{itemize}
  \item you know that \(f(x)\) is a polynomial of bounded degree,
  \item values \(f(0), f(1), \dots, f(n)\) are easy to compute,
  \item the asked point \(x\) is much larger than \(n\).
\end{itemize}

\section*{Core Idea}

For a degree-\(n\) polynomial, \(n+1\) sample points determine the polynomial uniquely. Lagrange interpolation writes
\[
f(x) = \sum_{i=0}^{n} f(i)\prod_{j \ne i} \frac{x-j}{i-j}.
\]

In modular arithmetic, the denominators are handled with modular inverses.

\section*{Key Insight}

In competitive programming, the target is usually evaluation, not symbolic reconstruction. That means you can use a fast
evaluation formula directly and avoid storing the full expanded polynomial.

\section*{Worked Problem}

\paragraph{Problem.}
Suppose a sequence is known to be generated by a polynomial of degree at most \(d\), and you have computed the first
\(d+1\) values by DP. Evaluate the sequence at a huge index \(N\).

\paragraph{Why interpolation fits.}
Once the degree bound is known, the early DP values are enough. Interpolation turns those \(d+1\) samples into the value
at \(N\) in roughly linear time in \(d\).

\section*{Correctness Intuition}

Two degree-\(d\) polynomials that agree on \(d+1\) distinct points must be identical. So the interpolated polynomial is
the same polynomial that generated the samples.

\section*{Complexity Analysis}

The straightforward modular Lagrange formula on points \(0,1,\dots,n\) can be evaluated in \(O(n)\) after prefix and
suffix product preprocessing.

\section*{Implementation}

The code assumes:

\begin{itemize}
  \item the modulus is prime,
  \item sample points are \(0,1,\dots,n\),
  \item \texttt{y[i] = f(i)}.
\end{itemize}

\section*{Common Pitfalls}

\begin{itemize}
  \item Interpolating when the sequence is not actually polynomial.
  \item Forgetting to handle the case where the query point is already among the sample indices.
  \item Using a composite modulus without checking inverse existence.
\end{itemize}

\section*{Variants / Extensions}

\begin{itemize}
  \item Interpolation on arbitrary \(x_i\).
  \item Newton-form interpolation.
  \item Recovering coefficients explicitly when the problem really asks for the polynomial itself.
\end{itemize}

\section*{Practice Problems}

\begin{itemize}
  \item Evaluate a polynomial sequence at a huge index after computing the first values by DP.
  \item Sum-of-powers style formulas modulo a prime.
  \item Sequence problems where finite differences reveal a small degree.
\end{itemize}

\section*{References}

\begin{itemize}
  \item \href{https://cp-algorithms.com/algebra/polynomial.html}{cp-algorithms: Operations on Polynomials and Series}.
  \item \href{https://github.com/kth-competitive-programming/kactl}{KACTL}.
  \item \href{https://codeforces.com/catalog}{Codeforces Catalog}.
\end{itemize}
